%===============================================================================
\section{Ethical considerations}
\label{section:ethical_considerations}

% Type of ethics
Only robot ethical considerations with humans as the ethical subjects and the ethics of good society (as is defined in Coeckelbergh book Robot Ethics\:\cite{coeckelbergh_robot_2022}) is considered in this paper.
% Responsibility
Coeckelbergh\:\cite{coeckelbergh_robot_2022} raises the ethical consideration of who is responsible for a robot, and if it is ethical as a developer to just abandon the robot after it has been created, leaving all responsibilities behind. Coeckelbergh\:\cite{coeckelbergh_robot_2022} goes on to discuss the ethical dilemma of what type of robots could be considered, or not considered, ethical to create.
In this project the responsibility of the robots, their use, and the creation of these robots fall on the stakeholder. The developers have no choice in whether it is ethical to develop this kind of robots (the project must be done in order to pass the course). Furthermore, the developers are unable to follow along with their creation after the course ends, and can thus not be held responsible for the robot. Finally, the developers are not in charge of whether the projects findings and research is published or not.
% Loosing jobs
The research could result in robots competing with professional football players, possibly resulting in reducing the human football scene and professional players loosing their jobs. Entertainment, especially sport, is a prominent aspect in today's society, so should the previous consideration play out, then it could lead to a transformation of the sports entertainment. This would in turn have some consequences, most likely including economical, societal, social and more. Similar considerations are brought up by Coeckelbergh\:\cite{coeckelbergh_robot_2022}, however this paper specifically applied those considerations to this project.
% Applications
Multi agent research, and robots in general, could be applied to military and surveillance sectors. This in turn raises questions of privacy and harm to humans, and who is to be responsible should that happen. Additionally, like Coeckelbergh\:\cite{coeckelbergh_robot_2022} asks the reader to ponder, in the case that the robots are employed for surveillance then what if the robot does not respect human dignity? Military applications require serious considerations as well, the disconnect that a robot creates could be seen as making it easier to kill and start conflicts since it is easier to lack sympathy for a target when one never has to interact or see them. This can be further extended to the view of some \ac{ai} as "black boxes", the inability to know the exact workings of an \ac{ai} and what it will do, even the developers might not fully understand them. If the robot \ac{ai} harms someone or something that is should not have (it was not its intended function to do so), then who is to blame?


%-------------------------------------------------------------------------------

\begin{comment}
Examine any ethical issues related to the potential harm your project could cause, particularly if you are developing robots or systems with significant impact. Consider safety, unintended consequences, and how your project might affect users or society. Outline measures to minimise risks and ensure your project is designed and implemented responsibly
\end{comment}

\begin{comment}
% Documentation
It is important that everything is openly documented and proper version control is available to ensure transparency, fair play and proof of the true developers. This will be provided by the projects public \href{https://github.com/DVA490-474-Project-Course}{GitHub}\footnote{https://github.com/DVA490-474-Project-Course} page where all code and documentation will be available.
% Sustainability
Proper consideration needs to be taken in regards to the materials used and the disposal of components in the event of replacement or damage. The project will be strict in its disposal and storage of materials and components, following the data sheets and regulations.
% Other fields
Additionally, care must be taken to the possible applications of the research in other fields like warfare and surveillance. The \acs{mit} License will be used to ensure transparency and thus reduce the risk of covert use of the research in unintended ways. 
% Safety and reliability
Finally, proper safety and reliability precautions and practices should be taken to minimise potential damage and harm that the robots can cause. This goes for both damage to robot opponents and disruption of the game, as well as potential harm to humans or the environment. 
\end{comment}

%===============================================================================